\documentclass[otchet]{SCWorks}
% Тип обучения (одно из значений):
%    bachelor   - бакалавриат (по умолчанию)
%    spec       - специальность
%    master     - магистратура
% Форма обучения (одно из значений):
%    och        - очное (по умолчанию)
%    zaoch      - заочное
% Тип работы (одно из значений):
%    coursework - курсовая работа (по умолчанию)
%    referat    - реферат
%  * otchet     - универсальный отчет
%  * nirjournal - журнал НИР
%  * digital    - итоговая работа для цифровой кафедры
%    diploma    - дипломная работа
%    pract      - отчет о научно-исследовательской работе
%    autoref    - автореферат выпускной работы
%    assignment - задание на выпускную квалификационную работу
%    review     - отзыв руководителя
%    critique   - рецензия на выпускную работу

% * Добавлены вручную. За вопросами к @mchernigin
\usepackage{preamble}

\begin{document}

% Кафедра (в родительном падеже)
\chair{математической кибернетики и компьютерных наук}

% Тема работы
\title{Анализ сложности префикс-функции, Z-функции и алгоритма Кнута-Морриса-Пратта}

% Курс
\course{2}

% Группа
\group{211}

% Факультет (в родительном падеже) (по умолчанию "факультета КНиИТ")
% \department{факультета КНиИТ}

% Специальность/направление код - наименование
 \napravlenie{02.03.02 "--- Фундаментальная информатика и информационные технологии}
% \napravlenie{02.03.01 "--- Математическое обеспечение и администрирование информационных систем}
% \napravlenie{09.03.01 "--- Информатика и вычислительная техника}
%\napravlenie{09.03.04 "--- Программная инженерия}
% \napravlenie{10.05.01 "--- Компьютерная безопасность}

% Для студентки. Для работы студента следующая команда не нужна.
% \studenttitle{студентки}

% Фамилия, имя, отчество в родительном падеже
\author{Кочергина Антона Алексеевича}

% Заведующий кафедрой 
\chtitle{доцент, к.\,ф.-м.\,н.}
\chname{С.\,В.\,Миронов}

% Руководитель ДПП ПП для цифровой кафедры (перекрывает заведующего кафедры)
% \chpretitle{
%     заведующий кафедрой математических основ информатики и олимпиадного\\
%     программирования на базе МАОУ <<Ф"=Т лицей №1>>
% }
% \chtitle{г. Саратов, к.\,ф.-м.\,н., доцент}
% \chname{Кондратова\, Ю.\,Н.}

% Научный руководитель (для реферата преподаватель проверяющий работу)
\satitle{доцент, к.\,ф.-м.\,н.} %должность, степень, звание
\saname{Ю.\,Н.\,Кондратова}

% Руководитель практики от организации (руководитель для цифровой кафедры)
\patitle{доцент, к.\,ф.-м.\,н.}
\paname{С.\,В.\,Миронов}

% Руководитель НИР
\nirtitle{доцент, к.\,п.\,н.} % степень, звание
\nirname{В.\,А.\,Векслер}

% Семестр (только для практики, для остальных типов работ не используется)
\term{2}

% Наименование практики (только для практики, для остальных типов работ не
% используется)
\practtype{учебная}

% Продолжительность практики (количество недель) (только для практики, для
% остальных типов работ не используется)
\duration{2}

% Даты начала и окончания практики (только для практики, для остальных типов
% работ не используется)
\practStart{01.07.2024}
\practFinish{13.01.2024}

% Год выполнения отчета
\date{2026}

\maketitle

% Включение нумерации рисунков, формул и таблиц по разделам (по умолчанию -
% нумерация сквозная) (допускается оба вида нумерации)
\secNumbering

\tableofcontents

% Раздел "Обозначения и сокращения". Может отсутствовать в работе
% \abbreviations
% \begin{description}
%     \item ... "--- ...
%     \item ... "--- ...
% \end{description}

% Раздел "Определения". Может отсутствовать в работе
% \definitions

% Раздел "Определения, обозначения и сокращения". Может отсутствовать в работе.
% Если присутствует, то заменяет собой разделы "Обозначения и сокращения" и
% "Определения"
% \defabbr

\section{Реализация алгоритмов}
\subsection{Z-функция}
\begin{Verbatim}[
  numbers=left,
  breaklines=true,
  breakanywhere=true
]
vector<int> z_function(string s){
    int n = (int)s.size();
    vector<int> z(n, 0);
    int l = 0, r = 0;

    for (int i = 1; i < n; i++){
        //если текущая позиция внутри уже найденного отрезка
        if (i <= r)
            z[i] = min(r - i + 1, z[i - l]);
        while (i + z[i] < n && s[z[i]] == s[i + z[i]])
            z[i]++;
        //если новый Z-отрезок длиннее текущего, обновляем границы
        if (i + z[i] - 1 > r){
            r = i + z[i] - 1;
            l = i;
        }
    }

    return z;
}
\end{Verbatim}

\subsection{Префикс-функция}

\begin{Verbatim}[
  numbers=left,
  breaklines=true,
  breakanywhere=true
]
vector<int> prefix_function(string s) {
    int n = (int)s.size();
    vector<int> p(n, 0);
    for (int i = 1; i < n; i++) {
        //длина нового префикса не может превышать p[i-1] + 11
        int c = p[i - 1];
        //уменьшаем c, пока не найдём позицию, где символ s[cur] совпадёт с s[i]
        while (s[i] != s[c] && c > 0)
            c = p[c - 1];
        // здесь либо s[i] == s[cur], либо cur == 0
        if (s[i] == s[c])
            p[i] = c + 1; //увеличиваем длину префикса на 1
    }
    return p;
}
\end{Verbatim}

\subsection{Алгоритм Кнута-Морриса-Пратта}

\begin{Verbatim}[
  numbers=left,
  breaklines=true,
  breakanywhere=true
]
vector<int> KMP(string s, string t){
    vector<int> A;
    if (s.empty()) return A;
    int k = 0;  //количество совпавших символов
    vector<int> pi = prefix_function(s);  //префикс-функция от образца S
    int n = s.size();
    int m = t.size();

    for (int i = 0; i < m; i++){
        //пока есть совпадение и текущий символ не совпадает
        while (k > 0 && t[i] != s[k]){
            k = pi[k - 1];
        }
        //если символы совпали
        if (t[i] == s[k]) {
            k++;
        }
        // если нашли полное вхождение
        if (k == n) {
            A.push_back(i - n + 1); //позиция начала вхождения
            k = pi[k - 1];
        }
    }

    return A;
}
\end{Verbatim}

\newpage

\section{Анализ сложности алгоритмов}

\subsection{Z-функция}



Функция z\_function получает строку s.

Создание вектора z размера n и его инициализация нулями выполняется за время $O(n)$, где n — длина строки s:
\begin{verbatim}
vector<int> z(n, 0);
\end{verbatim}

Инициализация переменных l = 0, r = 0 выполняется за $O(1)$.

Основной цикл выполняется для i от 1 до n-1, то есть $O(n)$ итераций:
\begin{verbatim}
for (int i = 1; i < n; i++) {
    if (i <= r)
        z[i] = min(r - i + 1, z[i - l]);
    while (i + z[i] < n && s[z[i]] == s[i + z[i]])
        z[i]++;
    if (i + z[i] - 1 > r) {
        r = i + z[i] - 1;
        l = i;
    }
}
\end{verbatim}

Внутри цикла:
\begin{itemize}
    \item Проверка условия i <= r и вычисление z[i] через min --- $O(1)$.
    \item Цикл while: Переменная z[i] увеличивается на каждой итерации этого while. 
    Суммарное количество увеличений z[i] по всем i не может превышать n, так как z[i] никогда не превышает n - i,
    а общая длина всех сравнений ограничена длиной строки. Каждое успешное сравнение символов сдвигает
    потенциальную правую границу r. Таким образом, суммарно по всем вызовам while получим $O(n)$.
    \item Обновление границ l и r --- $O(1)$.
\end{itemize}

Общая сложность алгоритма составит:

$T = O(1) + O(n) + O(1) + O(n) · (O(1) + O(1) + O(1) + O(1) + O(1)) + O(n) = O(n)$

\subsection{Префикс-функция}

Функция prefix\_function получает строку s.

Создание вектора p размера n и его инициализация нулями выполняется за время $O(n)$, где n — длина строки s:
\begin{verbatim}
vector<int> p(n, 0);
\end{verbatim}

Инициализация переменной n выполняется за $O(1)$.

Основной цикл выполняется для i от 1 до n-1, то есть $O(n)$ итераций:
\begin{verbatim}
for (int i = 1; i < n; i++) {
int c = p[i - 1];
while (s[i] != s[c] && c > 0)
c = p[c - 1];
if (s[i] == s[c])
p[i] = c + 1;
}
\end{verbatim}

Внутри цикла:
\begin{itemize}
\item Инициализация переменной c = p[i - 1] — $O(1)$.
\item Цикл while: переменная c уменьшается на каждой итерации этого while.
Значение p[i] может увеличиться не более чем на 1 по сравнению с p[i-1], а всего за все итерации внешнего цикла
p[i] может возрасти не более чем на n. Каждое уменьшение c в цикле while соответствует хотя бы одному
предыдущему увеличению, поэтому суммарное количество итераций while по всем i не превышает n.
Таким образом, суммарно по всем вызовам while получим $O(n)$.
\item Проверка условия s[i] == s[c] и присваивание p[i] = c + 1 — $O(1)$.
\end{itemize}

Общая сложность алгоритма составит:

$T = O(1) + O(n) + O(1) + O(n) \cdot (O(1) + O(1) + O(1)) + O(n) = O(n)$

\subsection{Алгоритм Кнута-Морриса-Пратта}

Функция KMP получает строки s (образец) и t (текст).


Инициализация переменной k = 0 выполняется за $O(1)$.

Вычисление префикс-функции для строки s (образца) выполняется за время $O(n)$, где n — длина строки s:
\begin{verbatim}
vector<int> pi = prefix_function(s);
\end{verbatim}

Инициализация переменных n = s.size() и m = t.size() выполняется за $O(1)$.

Основной цикл выполняется для i от 0 до m-1, то есть $O(m)$ итераций, где m — длина строки t:
\begin{verbatim}
for (int i = 0; i < m; i++) {
while (k > 0 && t[i] != s[k]) {
k = pi[k - 1];
}
if (t[i] == s[k]) {
k++;
}
if (k == n) {
A.push_back(i - n + 1);
k = pi[k - 1];
}
}
\end{verbatim}

Внутри цикла:
\begin{itemize}
\item Цикл while: переменная k уменьшается на каждой итерации этого while.
Значение k может увеличиться не более чем на 1 за одну итерацию внешнего цикла (при совпадении символов),
а всего за все итерации внешнего цикла k может возрасти не более чем на m. Каждое уменьшение k в цикле while
соответствует хотя бы одному предыдущему увеличению, поэтому суммарное количество итераций while по всем i
не превышает m. Таким образом, суммарно по всем вызовам while получим $O(m)$.
\item Проверка условия t[i] == s[k] и увеличение k — $O(1)$.
\item Проверка условия k == n и добавление индекса в вектор A — $O(1)$.
\item Обновление k = pi[k - 1] при нахождении вхождения — $O(1)$.
\end{itemize}

Общая сложность алгоритма составит:

$T = O(1) + O(1) + O(n) + O(1) + O(m) \cdot (O(1) + O(1) + O(1)) + O(m) = O(n + m)$

\end{document}
